\section{Introduction}\label{sec:intro}

Haiden, Katzarkov, Kontsevich and Pandit \cite{hkkp26} propose a notion of
K\"ahler geometry on a triangulated category with a stability condition in
the sense of Bridgeland \cite{bri07}. In their framework an object carries a
space of metrics, a complex-valued function $\cS_{\CC}$ on it whose
differential is given by an endomorphism $\Omega$, a phase $\arg\Omega$
defined by a positivity condition, a mass $\tr|\Omega|$, and a flow
$h^{-1}\dot h=-\arg\Omega$ of the metrics, called the minimizing flow in
\cite[\S4.1]{hkkp26}; a metric is \emph{harmonic} if $\arg\Omega$ is a
scalar. For a representation of a quiver with central charges $z_i$ and
Hermitian metrics, $\Omega$ is $z_i$ shifted by the moment map at each
vertex, and the function $\cS_{\CC}$ is explicit \cite[\S5.1]{hkkp26}: after
rotation by the argument of the central charge of the representation, its
imaginary part is a positive multiple of the Kempf--Ness function. For a
complex of representations an endomorphism $\Omega$ is defined in
\cite[\S5.4]{hkkp26} through a superconnection in a graded algebra attached
to the quiver; its trace is the central charge of the complex. No function
$\cS_{\CC}$ is given there for complexes.

For a representation the flow is not the gradient flow of the Kempf--Ness
function, but the Kempf--Ness function does not increase along it
(Proposition~\ref{prop:kempfness}); that is, it is a Lyapunov function for
the flow. The question of this paper is whether the Kempf--Ness function
of a representation extends to complexes. Two properties of the Kempf--Ness
function \cite{kn79,kin94} are meant: its differential is the moment map
shifted by the charges, and it does not increase along the flow. As a
function with differential $\Omega$ it extends to every bounded complex of
representations of every finite quiver (Theorem~\ref{thm:introA}). As a
function that does not increase along the flow it extends to two classes of
complexes (Theorem~\ref{thm:introC}) and not to all
(Theorem~\ref{thm:introD}). Even on a complex with polystable cohomology, the phase potential can be
unbounded below and strictly concave along a geodesic in the space of
metrics (Corollary~\ref{cor:increase}).

In \cite[\S5.4]{hkkp26} the phase of $\Omega$ is written with the functional
calculus, which presupposes that $\Omega$ is normal. For complexes this can
fail (Example~\ref{ex:nonnormal}). In the general framework of
\cite[\S4.2]{hkkp26} the argument is a self-adjoint element $\arg\Omega$ such
that $\Omega(e^{-\sqrt{-1}\arg\Omega}x^*x)\geq0$ for all elements $x$ of the
algebra attached to the metrized object, $\Omega$ being a functional there;
for the functional $x\mapsto\tr(\Omega x)$ on endomorphisms this says that
$\Omega e^{-\ii\arg\Omega}$ is positive semidefinite. We therefore define
$\arg\Omega$ as the Hermitian endomorphism $\Phi$, with spectrum in the
interval $(-\pi n,\pi(1-n))$ in degree $n$, the branch of
\cite[(5.15)]{hkkp26}, for which $\Omega e^{-\ii\Phi}$ is positive definite;
$e^{\ii\Phi}$ is the unitary factor of the polar decomposition of $\Omega$.
For normal $\Omega$ this is the argument of the functional calculus. All
statements about the flow of a complex refer to this definition. We call the
resulting flow the \emph{phase flow}; we do not use the name minimizing flow
because on complexes the mass can increase along it
\cite[Theorem~1.1]{wue26a}. It is the flow of \cite{wue26b,wue26c}, whose
titles keep the name of \cite{hkkp26}.

The phase flow was studied in \cite{wue26a,wue26b} for two-term complexes of
representations of a generalized Kronecker quiver that are quasi-isomorphic
to a representation: the acyclic part decays, and for suitable charges the
flow converges on the cohomology when the cohomology is polystable
\cite[Theorems~1.1 and~1.2]{wue26b}; these two notions are defined in
Section~\ref{sec:setting}. For these complexes a function with
differential $\Omega$ was found by direct computation
\cite[Proposition~7.1]{wue26a}, and for a range of charges its rotated
imaginary part decreases along the flow \cite[Theorem~7.6]{wue26b}.

\subsection{Results}
Let $Q$ be a finite quiver with central charges $z_i$ in the upper
half-plane, and $\mathsf E$ a bounded complex of representations of $Q$ with
spaces $E_i^n$ at the vertex $i$ in degree $n$, differentials $d_i$, arrows
$\rho(a)$ and a Hermitian metric $h_i^n$ on each $E_i^n$
(Section~\ref{sec:setting}). The endomorphism $\Omega$ preserves each
$E_i^n$; we write $\Omega_i^n$ for its restriction to $E_i^n$, and
$\norm\cdot_{\HS}$ for the Hilbert--Schmidt norm of a linear map for the
metrics. An integral formula for $\Omega$ is given in
Proposition~\ref{prop:coefficient}. Its trace is the central charge
$Z=\sum_{i,n}(-1)^nz_i\dim E_i^n$ of the complex.

\begin{theorem}[Primitive]\label{thm:introA}
For every bounded complex of representations of $Q$ with a metric:
\begin{enumerate}
\item[(i)] The one-form $\xi\mapsto\tr(\Omega\,\xi)$ on the space of metrics,
 $\xi=h^{-1}\delta h$, is closed.
\item[(ii)] For $x\in[0,1]$ let $\Omega^{(x)}$ be the endomorphism of the
 complex with differentials $\sqrt x\,d_i$ and the same arrows and metrics.
 Then
 \[
 \cS_{\CC}(h)=\sum_n(-1)^n\Bigl[\sum_iz_i\log\det h_i^n-\sum_a\norm{\rho^n(a)}^2_{\HS}\Bigr]
 +\int_0^1\sum_{i,n}n\,\tr\bigl((\Omega^{(x)})_i^n-(\Omega^{(0)})_i^n\bigr)\,\frac{dx}x
 \]
 satisfies $d\cS_{\CC}(\xi)=\tr(\Omega\,\xi)$. Here $\rho^n(a)$ is the
 restriction of $\rho(a)$ to degree $n$, and the determinants are Gram
 determinants in fixed bases.
\item[(iii)] If $Z$ is not zero, then for any real $\varphi$ with
 $Z=|Z|e^{\ii\varphi}$, along the phase flow
 \[
 \frac d{dt}\Ima\bigl(e^{-\ii\varphi}\cS_{\CC}\bigr)
 =-\sum_{i,n}\tr\Bigl(|\Omega_i^n|\,(\Phi_i^n-\varphi)\sin(\Phi_i^n-\varphi)\Bigr),
 \]
 where $\Omega_i^n=|\Omega_i^n|e^{\ii\Phi_i^n}$ is the polar decomposition
 with the spectrum of $\Phi_i^n$ in the interval
 \mbox{$(-\pi n,\pi(1-n))$}.
\end{enumerate}
\end{theorem}

We call $\Ima(e^{-\ii\varphi}\cS_{\CC})$ the \emph{phase potential}. For a
representation in degree zero, with $\varphi\in(0,\pi)$, it is
$\sin\varphi$ times the Kempf--Ness function
$\sum_a\norm{\rho(a)}^2_{\HS}+\sum_i\theta_i\log\det h_i$ with
$\theta_i=\Ima(e^{-\ii\varphi}z_i)/\sin\varphi$, and each trace in the sum
in (iii) is nonnegative (Proposition~\ref{prop:kempfness}). For complexes
they can have either sign. The last term of $\cS_{\CC}$ is a transgression by
the number operator along the rescaling of the differentials, of the kind
used for Bott--Chern forms of Hermitian complexes \cite{bgs88}, here taken on
the space of metrics and in the algebra attached to the quiver.

Let $K_m$ be the generalized Kronecker quiver, with vertices $0,1$ and $m$
arrows from $0$ to $1$. We write a complex of representations of $K_m$ as a
row of its terms, each term being a representation written as (space at the
vertex $0$ $\to$ space at the vertex $1$); the maps between consecutive terms
are the differentials.

\begin{theorem}[Monotonicity in two classes]\label{thm:introC}
In each of the following cases the phase potential is nonincreasing along the
phase flow.
\begin{enumerate}
\item[(i)] $i_0$ is a vertex of $Q$ without loops; the complex is a vector
 space $E_i^0$ in degree zero at every vertex $i\neq i_0$ and $V\to U$ in
 degrees $0,1$ at $i_0$, with surjective differential and nonzero
 cohomology; and, for $\varphi=\arg Z\in(0,\pi)$, one has
 $\arg z_{i_0}\geq\varphi$ and either $\varphi\geq\frac{\pi}3$ or
 $\dim U=1$.
\item[(ii)] $Q=K_m$ and $\arg z_0\geq\arg z_1$; the complex is
 $(0\to U_1)\to(V\to W)\to(U_0\to0)$ in degrees $-1,0,1$, that is, with $U_1$
 at the vertex $1$ in degree $-1$ and $U_0$ at the vertex $0$ in degree $1$
 (the subscript is the vertex), with injective differential $U_1\to W$,
 surjective differential $V\to U_0$ and nonzero cohomology; and
 $\dim U_0\leq1$, $\dim U_1\leq1$.
\end{enumerate}
\end{theorem}

\begin{theorem}[Failure of monotonicity]\label{thm:introD}
Let $Q=K_m$, $z_0=-1+\ii$ and $z_1=1+\ii$.
\begin{enumerate}
\item[(i)] Let $\mathsf E$ be the orthogonal direct sum of a representation
 with spaces of equal dimension $k\geq1$ at the two vertices, with a harmonic
 metric, and of the complex $\CC\to\CC$ at the vertex $1$ in degrees $0,1$
 with a nonzero differential $f$. Then the phase potential increases along
 the whole trajectory, at the rate $\pi e^{-|f|^2}$, where $|f|$ is the norm
 of $f$ for the metrics at the time considered, and it is not bounded below
 on the metrics of $\mathsf E$.
\item[(ii)] There is a three-term complex as in
 Theorem~\ref{thm:introC}(ii) with $\dim U_0=2$, $\dim U_1=1$ and stable
 cohomology, with a metric at which the phase potential increases along the
 flow.
\end{enumerate}
\end{theorem}

Theorem~\ref{thm:introC} is proved as Theorems~\ref{thm:onevertex}
and~\ref{thm:rankone}, and Theorem~\ref{thm:introD} as
Corollary~\ref{cor:increase} and Theorem~\ref{thm:twosidedfailure}.

\emph{Sharpness.} For $Q=K_m$, $i_0=0$ and cohomology at both vertices,
$\varphi$ lies between $\arg z_1$ and $\arg z_0$, so the condition
$\arg z_0\geq\varphi$ is $\arg z_0\geq\arg z_1$. For $\arg z_0>\arg z_1$,
strict decrease when $\varphi\geq 2\pi/5$ or $\dim U=1$ is
\cite[Theorem~7.6]{wue26b}. The stronger bound $\pi/3$ follows from the elementary trace estimate
in Section~\ref{sec:setting}. Theorem~\ref{thm:introC}(i) adds
arbitrary quivers, the case of equality, and cohomology at a single vertex. In Theorem~\ref{thm:introD}(i) the differential is at
one vertex and $\arg z_1<\varphi$; so the condition
$\arg z_{i_0}\geq\varphi$ of Theorem~\ref{thm:introC}(i) cannot be dropped.
For $\dim U\geq2$ a lower bound on $\varphi$ cannot be dispensed with
either: for $K_2$, the charges $z_0=10^{-6}(-1+\ii)$, $z_1=1+\frac\ii{500}$
and cohomology of dimension vector $(1,1)$, where
$\varphi=0.0020\ldots<\frac{\pi}3$, there is a complex with $\dim U=2$ and a
metric at which the potential increases \cite[Proposition~10.2]{wue26b}. The optimal lower bound on $\varphi$ for
nonnegative total dissipation at every metric remains open. Theorem~\ref{thm:introD}(ii) shows that the
condition on the dimensions in Theorem~\ref{thm:introC}(ii) cannot be
dropped; for $z_0=-1+\ii$, $z_1=1+\ii$ and cohomology with equal dimensions
at the two vertices, the pairs $(\dim U_0,\dim U_1)$ for which the phase
potential is nonincreasing at all metrics are determined in
Corollary~\ref{cor:classification}.

\emph{Method.} The closedness in Theorem~\ref{thm:introA} needs only two odd
derivations of a graded algebra with the relations of
Lemma~\ref{lem:abstract}, a functional that vanishes on
supercommutators and on the images of the derivations, and a formula for the
variation of the curvature element under a change of metric; the same lemma
serves the algebra of the quiver
and, for bundles, the algebra of endomorphism-valued forms, where it is the
classical argument by which Bott--Chern forms do not depend, modulo
$\partial$- and $\bar\partial$-exact forms, on the path of metrics. Theorem~\ref{thm:introC}(i) uses the trace estimate
proved after Theorem~\ref{thm:imported}, which strengthens
\cite[Theorem~3.1]{wue26b}. The entries of the Schur multiplier are the
second divided differences of $e^{-x}$ at $\ell_i,0,\ell_k$, where the
$\ell_i$ are the eigenvalues of the differential times its adjoint. In the three-term complexes the point $0$ is
replaced by the eigenvalues of the other differential times its adjoint, and
no such inequality holds uniformly in these eigenvalues
(Remark~\ref{rem:interpolation}); this is the source of the example in
Theorem~\ref{thm:introD}(ii). The numerical values in that example are
enclosures obtained with interval arithmetic \cite{joh17}.

\begin{theorem}[Duality]\label{thm:introB}
Let $\mathsf E^\vee$ be the dual complex over the opposite quiver, with
$(E^\vee)_i^n=(E_i^{-n})^*$, the transposed maps, the dual metrics and the
charges $-\overline{z_i}$. Then $\Omega^\vee=-(\Omega^\dagger)^{\mathsf T}$,
the phases satisfy $\Phi^\vee=\pi-\Phi^{\mathsf T}$, the masses agree, and
$h\mapsto e^{-\pi t}(h^{-1})^{\mathsf T}$ maps trajectories of the phase flow
of $\mathsf E$ to trajectories of the phase flow of $\mathsf E^\vee$. If
$Z(\mathsf E)\neq0$ and the argument of $Z(\mathsf E^\vee)$ is taken as
$\pi-\varphi$, this map preserves the phase potential up to a constant.
\end{theorem}

Duality is used as a tool: it halves the cases in
Theorem~\ref{thm:onevertex}, and it carries statements about complexes with a
surjective differential in degrees $0,1$ over to complexes with an injective
differential in degrees $-1,0$ (Remark~\ref{rem:injective}).

\subsection{Bundles}
The last section gives the corresponding formulation for a two-term complex
$E^0\to E^1$ of holomorphic vector bundles with Hermitian metrics on a
compact Riemann surface, with $-\frac{\ii}{2\pi}\Lambda F_h$ in the place of
the moment map, following the dictionary of \cite[Remark~5.14]{hkkp26} and
using Quillen's superconnection \cite{qui85}. We record the formal
properties, which for the closedness follow the theory of Bott--Chern forms
\cite{bc65,don85,bgs88}: the one-form defined by $\Omega$ on the space of
metrics is closed, with a primitive given by the same transgression; for a
single bundle the phase potential is a positive multiple of Donaldson's
functional \cite{don85}; the derivative formula of
Theorem~\ref{thm:introA}(iii) holds along smooth solutions; the
linearization is parabolic in the sense of Petrovskii
for the convention $\Delta_\omega=-d^*d$; and for a section $s$ of a
line bundle the sum $(\Omega^0+\Omega^1)\,\omega$ of the two blocks of
$\Omega$, times the K\"ahler form, is a representative of the Euler class of
the kind constructed by Mathai and Quillen \cite{mq86}, and the integrand of
the phase potential decreases at every point of the surface
(Propositions~\ref{prop:bundleprimitive}--\ref{prop:section}).

\subsection{Open questions}
Convergence of the phase flow on the cohomology and decay of the acyclic
part are not treated here, except for the contractible summands of
Proposition~\ref{prop:contractible}. For two-term complexes over $K_m$ with
a surjective differential at the vertex $0$, the decay is proved in
\cite{wue26b} for all charges and the convergence for strictly ordered
charges and polystable cohomology. For the corresponding complexes on
acyclic quivers, \cite{wue26c} proves decay of the exact part and
exponential convergence of the rescaled metrics on polystable
Harder--Narasimhan quotients of the cohomology.
Remark~\ref{rem:injective} carries these results to injective differentials. For the
three-term complexes of Theorem~\ref{thm:introC}(ii) both are open. Beyond
Theorems~\ref{thm:onevertex} and~\ref{thm:rankone} and
Corollaries~\ref{cor:increase} and~\ref{cor:classification} we do not know
when the phase potential is monotone. For bundles, existence of solutions of
the flow is not proved, and neither is the convergence along the flow, for a
section of a line bundle, of the form $(\Omega^0+\Omega^1)\,\omega$ to the
current of integration over the zeros of the section.

\subsection{Organization}
Section~\ref{sec:setting} recalls the definition of \cite{hkkp26}, gives the
integral formula for $\Omega$, fixes the terminology, and sharpens one of the two
matrix inequalities recalled from \cite{wue26b}. Section~\ref{sec:primitive}
proves Theorem~\ref{thm:introA}(i),(ii), and Section~\ref{sec:rate} proves
(iii) and Theorem~\ref{thm:introD}(i). Section~\ref{sec:duality} proves
Theorem~\ref{thm:introB}. Section~\ref{sec:monotone} proves
Theorems~\ref{thm:introC} and~\ref{thm:introD}(ii).
Section~\ref{sec:bundles} is on bundles.
